\section{Details on the Receiver-Relative Rate-Distortion View}
\label{app:rate-distortion}

This appendix provides additional details for the local arguments used in
Section~\ref{sec:information_theory}.
The purpose of these derivations is not to establish an exact rate-distortion
theory of LoRA, but to make explicit the assumptions under which the geometric
quantities introduced in the main text arise.

\subsection{Local Behavioral Distortion}
\label{app:local-distortion}

Let $A$ denote a trained adapter and let $\hat A=A+e$ be a perturbed version,
where $e$ is the error induced by compression.
For a differentiable evaluation loss $\mathcal{L}_U$, a Taylor expansion around
$A$ gives
\begin{equation}
\begin{split}
    \mathcal{L}_U(M\oplus(A+e))
    -
    \mathcal{L}_U(M\oplus A)
    &=
    g_U^\top e
    +
    \frac{1}{2}e^\top H_Ue
    +
    o(\|e\|^2),
\end{split}
\label{eq:app-taylor}
\end{equation}
where
\[
    g_U
    =
    \nabla_A \mathcal{L}_U(M\oplus A),
    \qquad
    H_U
    =
    \nabla_A^2 \mathcal{L}_U(M\oplus A).
\]

The quadratic approximation used in the main text therefore requires the
first-order contribution to be negligible for the perturbations considered.
This is exact at a stationary point of the evaluation objective and can also
be a useful local approximation when $|g_U^\top e|$ is small compared with the
quadratic term.

The exact Hessian need not be positive semidefinite.
When interpreting $H_U$ as a local distortion metric, one may instead use a
positive-semidefinite curvature approximation, such as a generalized
Gauss--Newton matrix.
The relevant point for our argument is not the particular approximation, but
that behavioral distortion generally defines a non-isotropic metric over
parameter perturbations.

\subsection{When Does Weight Error Recover Functional Distortion?}
\label{app:isotropy}

Squared weight reconstruction error corresponds to the Euclidean quadratic
form
\[
    D_{\mathrm{weight}}(e)=e^\top e,
\]
whereas the local functional criterion is
\[
    D_{\mathrm{func}}(e)=e^\top H_Ue.
\]

Let $\mathcal{S}$ denote the subspace containing the compression errors under
consideration.
If
\begin{equation}
    H_U|_{\mathcal{S}}=\lambda I,
    \qquad \lambda>0,
    \label{eq:isotropic-condition}
\end{equation}
then
\[
    D_{\mathrm{func}}(e)
    =
    \lambda D_{\mathrm{weight}}(e)
\]
for every $e\in\mathcal{S}$, and the two criteria induce identical rankings of
compression errors.

Conversely, if all perturbations with equal Euclidean norm have equal
functional cost, the restriction of $H_U$ to $\mathcal{S}$ must be proportional
to the identity.
To see this, diagonalize the symmetric quadratic form on $\mathcal{S}$.
If two eigenvalues differ, perturbations of equal norm along their associated
eigenvectors have different functional costs.

Thus, ordinary reconstruction error is sufficient only under an isotropic
functional geometry on the errors produced by the codec.
The RMSE counterexample and layer-location controls in the main text provide
empirical evidence against this approximation in our setting.

\subsection{Correction Statistics and the Learned Update}
\label{app:correction-statistics}

At the frozen parameters $\theta_0$, let
\[
    g_i=\nabla_\theta\ell_i(\theta_0),
    \qquad
    \mu=\mathbb{E}[g_i],
    \qquad
    J=\mathbb{E}[g_ig_i^\top].
\]
The centered correction covariance is
\begin{equation}
    C_g
    =
    \mathbb{E}\big[
        (g_i-\mu)(g_i-\mu)^\top
    \big]
    =
    J-\mu\mu^\top.
    \label{eq:app-centered-gradient}
\end{equation}

A quadratic approximation to the average training objective has the form
\[
    \mathcal{L}_T(\theta_0+\Delta\theta)
    \approx
    \mathcal{L}_T(\theta_0)
    +
    \mu^\top\Delta\theta
    +
    \frac{1}{2}
    \Delta\theta^\top H_T\Delta\theta.
\]
When $H_T$ is invertible on the relevant subspace, its minimizer is
\begin{equation}
    \Delta\theta^\star
    =
    -H_T^{-1}\mu.
    \label{eq:app-newton-update}
\end{equation}

Throughout, $H_T^{-1}$ denotes the inverse restricted to the relevant
subspace. When $H_T$ is singular, the same local expressions can instead be
read using a pseudoinverse or a suitably damped inverse.

Equation~\ref{eq:app-newton-update} makes clear that $J$ does not determine the
deterministic local update.
For example, in one dimension, a dataset for which every example has gradient
$+1$ and a dataset containing equal numbers of gradients $+1$ and $-1$ both
satisfy
\[
    J=1,
\]
but their means are respectively $\mu=1$ and $\mu=0$.
Their local optimal updates are therefore different despite having identical
second moments.

We consequently interpret $J$ as a descriptor of the geometry of the
corrections requested by the data, rather than as an estimate of the adapter
update itself.
Its potential relevance to compression comes from the directional diversity
that it captures in addition to the mean correction.

\subsection{From Correction Statistics to a Receiver-Relative Spectrum}
\label{app:relative-spectrum}

To motivate a link between correction statistics and compressibility, consider
the following local linear model.
Let $G$ be a random correction gradient and suppose that it induces a local
response
\begin{equation}
    \Delta\theta=-H_T^{-1}G.
    \label{eq:app-linear-response}
\end{equation}
This should be understood as an idealized response model, not as a description
of an SGD trajectory.

Under Equation~\ref{eq:app-linear-response}, the mean update is
\[
    \mathbb{E}[\Delta\theta]
    =
    -H_T^{-1}\mu,
\]
its covariance is
\begin{equation}
    \mathrm{Cov}(\Delta\theta)
    =
    H_T^{-1}C_gH_T^{-1},
    \label{eq:app-update-covariance}
\end{equation}
and its second moment is
\begin{equation}
    \mathbb{E}[
        \Delta\theta\Delta\theta^\top
    ]
    =
    H_T^{-1}JH_T^{-1}.
    \label{eq:app-update-second-moment}
\end{equation}

Behavioral distortion introduces a second geometry through $H_U$.
Applying it to the centered update covariance in
Equation~\ref{eq:app-update-covariance} gives
\begin{equation}
    C_{\mathrm{rel}}
    =
    H_U^{1/2}
    H_T^{-1}
    C_g
    H_T^{-1}
    H_U^{1/2}.
    \label{eq:app-relative-covariance}
\end{equation}
This is the quantity that corresponds most directly to the Gaussian
rate-distortion argument below: it is the covariance of the local update
distribution expressed in behaviorally weighted coordinates.

For prediction, however, we are also interested in the systematic correction
encoded by $\mu$. Since
\[
    J=C_g+\mu\mu^\top,
\]
retaining the uncentered second moment motivates the related geometry
\begin{equation}
    \widetilde C_{\mathrm{rel}}
    =
    H_U^{1/2}
    H_T^{-1}
    J
    H_T^{-1}
    H_U^{1/2}.
    \label{eq:app-relative-second-moment}
\end{equation}
Unlike $C_{\mathrm{rel}}$, $\widetilde C_{\mathrm{rel}}$ is not itself a
Gaussian source covariance. It is a second-moment extension that retains both
the mean correction and its directional variability, and is introduced as a
candidate geometry for predicting behavioral rate.

When training and evaluation share the same local geometry,
$H_T=H_U=H$, this reduces to
\begin{equation}
    \widetilde C_{\mathrm{rel}}
    =
    H^{-1/2}JH^{-1/2}.
    \label{eq:app-whitened-spectrum}
\end{equation}

This whitening has a simple interpretation.
A correction direction is discounted when the training geometry makes it easy
to realize, and amplified when it corresponds to a behaviorally sensitive
direction of the receiver.

Our empirical implementation is more limited than Equation~\ref{eq:app-relative-second-moment}: the $J$ used later is estimated from projected output-head correction gradients rather than gradients in the full LoRA parameter space.
It is therefore a correction-gradient second moment, not the Fisher information matrix, the Hessian, or the learned update itself.
We use it as an accessible proxy for the receiver-relative geometry motivated above.

\subsection{Gaussian Rate-Distortion Motivation}
\label{app:gaussian-rd}

The spectral summaries used in the main text are motivated by the classical
rate-distortion function of a Gaussian source under quadratic distortion
\citep[Ch.~10]{cover2006elements}.

Let
\[
    X\sim\mathcal{N}(0,\Sigma)
\]
and consider a quadratic distortion
\[
    d(x,\hat x)
    =
    (x-\hat x)^\top H(x-\hat x),
    \qquad H\succeq0.
\]
After the change of coordinates
\[
    Z=H^{1/2}X,
\]
the problem becomes ordinary squared-error coding for a Gaussian variable with
covariance
\[
    C=H^{1/2}\Sigma H^{1/2}.
\]
Let $\lambda_1,\ldots,\lambda_d$ denote the eigenvalues of $C$.
The Gaussian rate-distortion function is then given by reverse water filling:
\begin{equation}
    R(D)
    =
    \frac{1}{2}
    \sum_j
    \left[
        \log_2\frac{\lambda_j}{\nu}
    \right]_+,
    \label{eq:app-waterfilling}
\end{equation}
where $\nu$ is chosen so that
\[
    D
    =
    \sum_j\min(\lambda_j,\nu).
\]

Equation~\ref{eq:app-waterfilling} highlights two properties relevant to our
setting.
Rate increases when the distortion-weighted variance becomes larger, but also
when it is spread across more active directions.
This motivates considering both scale-sensitive and
dimensionality-sensitive summaries of correction spectra.

A convenient smooth summary is
\begin{equation}
    I_{\mathrm{rel}}(\alpha)
    =
    \frac{1}{2}
    \log\det
    \left(
        I+\alpha C_{\mathrm{rel}}
    \right)
    =
    \frac{1}{2}
    \sum_j
    \log(1+\alpha\lambda_j).
    \label{eq:app-soft-logdet}
\end{equation}
Unlike Equation~\ref{eq:app-waterfilling}, this expression does not implement
an explicit distortion threshold and should not be interpreted as the
rate-distortion function of our adapters.
We use it only as a smooth statistic that increases with both the magnitude
and dimensionality of the relevant spectrum.

Our empirical correction spectrum is a substantially coarser approximation.
It is constructed from projected output-head correction gradients rather than
from gradients in the full LoRA parameter space, and we do not explicitly
estimate $H_T$ or $H_U$ in the experiments reported here.
Consequently,
\[
    \log\det(I+\alpha J)
\]
should be interpreted as a tractable proxy for the second-moment geometry
$\widetilde C_{\mathrm{rel}}$, which is itself an extension of the covariance
geometry directly motivated by Gaussian rate-distortion theory.
It is therefore neither a direct estimate of
Equation~\ref{eq:app-relative-covariance} nor an information-theoretic rate.